\abstract{Digital cognitive training is often delivered through stationary screen interaction, whereas everyday self-regulation occurs while the child moves through a variable physical environment. This companion study tests whether embodied action and augmented reality (AR) contribute to contextual generalization of the ARC-ADHD challenge-regulation-recovery rule. The design explicitly separates three conditions: stationary touchscreen interaction, embodied movement without AR, and embodied AR interaction. This decomposition prevents movement, novelty, camera use, spatial mapping, and augmented content from being treated as a single mechanism. Matched tasks preserve the cognitive rule while changing the action and environmental context. Primary proximal outcomes are recovery latency and cognitive-control performance; movement exposure, sensorimotor synchronization, spatial behavior, tolerability, and safety are secondary mechanistic outcomes. Cognitive performance in embodied and AR conditions is evaluated with a preregistered non-inferiority margin because increased movement is not considered beneficial if it degrades control. A later longitudinal test can examine whether repeated contextual variation improves independent transfer. The study treats AR as an experimental context rather than as a therapeutic mechanism by definition.}
